% Diagrama da arquitetura do SIM — incluído via \input
\begin{figure}[htbp]
    \centering
    \resizebox{0.92\linewidth}{!}{%
    \begin{tikzpicture}[
        font=\small,
        >={Stealth[length=2mm]},
        box/.style={draw, rounded corners=2pt, align=center, minimum width=2.5cm,
                    minimum height=1.0cm, fill=white, line width=0.5pt},
        store/.style={box, fill=black!7},
        flow/.style={->, thick, line width=0.6pt},
        biflow/.style={<->, thick, line width=0.6pt},
        elabel/.style={font=\footnotesize, fill=white, inner sep=1.5pt},
        zone/.style={rounded corners=5pt, draw=black!30, dash pattern=on 2pt off 2pt,
                     inner sep=0.3cm},
        zonelbl/.style={font=\scriptsize\bfseries, color=black!55, inner sep=2pt},
        node distance=1.05cm and 0.95cm
    ]
        % ===== Camada 1: canal conversacional (produção) =====
        \node[box, minimum width=2.6cm] (user) {Usuário\\\textit{WhatsApp}};
        \node[box, right=of user] (waha) {Gateway\\WAHA};
        \node[box, right=of waha]  (n8n)  {Automação\\n8n};
        \node[box, right=of n8n, minimum width=2.8cm] (orch)
              {Agente orquestrador\\[1pt]{\scriptsize\texttt{GPT-4o}}};

        \draw[flow] (user) -- (waha);
        \draw[flow] (waha) -- (n8n);
        \draw[flow] (n8n)  -- (orch);

        % retorno: resposta em linguagem natural (roteada acima da linha)
        \coordinate (top) at ($(orch.north)+(0,0.9)$);
        \draw[flow] (orch.north) -- (top) -| (user.north);
        \node[elabel] at ($(user.north)!0.5!(orch.north)+(0,0.9)$)
              {resposta em linguagem natural};

        % ===== Camada 2: recuperação e consulta =====
        \node[box, below=1.8cm of n8n,  minimum width=3.0cm] (rag)
              {RAG documental\\{\scriptsize Diário Oficial}};
        \node[box, below=1.8cm of orch, minimum width=3.0cm] (sql)
              {Text-to-SQL\\[1pt]{\scriptsize\texttt{GPT-4.1-mini}}};

        \coordinate (sp) at ($(orch.south)+(0,-0.75)$);
        \draw[flow] (orch.south) -- (sp);
        \draw[biflow] (sp) -| (rag.north);
        \draw[biflow] (sp) -- (sql.north);
        \node[elabel] at ([yshift=0.34cm]rag.north) {normativo};
        \node[elabel] at ([yshift=0.34cm]sql.north) {tabular / analítico};

        % ===== Camada 3: dados =====
        \node[store, below=0.9cm of rag] (qdrant)
              {Índice vetorial\\Qdrant};
        \node[store, below=0.9cm of sql] (duck)
              {API analítica\\DuckDB \textit{in-memory}};
        
        \node[store, minimum width=3.8cm, below=0.9cm of qdrant] (ingest)
              {Indexação Documental\\{\scriptsize Diário Oficial (PDFs / OCR)}};
        \node[store, minimum width=3.8cm, below=0.9cm of duck] (etl)
              {ETL SICOM e CGU\\{\scriptsize DDL e Tabelas Consolidadas}};

        \draw[biflow] (rag) -- (qdrant);
        \draw[biflow] (sql) -- (duck);
        \draw[flow] (ingest) -- (qdrant);
        \draw[flow] (etl) -- (duck);

        % ===== Zonas de fundo =====
        \begin{scope}[on background layer]
            \node[zone, fit=(user)(waha)(n8n)(orch)] (zA) {};
            \node[zone, fit=(rag)(sql)]               (zB) {};
            \node[zone, fit=(qdrant)(duck)(ingest)(etl)] (zC) {};
        \end{scope}
        \node[zonelbl, anchor=south west] at (zA.north west) {Canal conversacional (produção)};
        \node[zonelbl, anchor=south west] at (zB.north west) {Recuperação e consulta};
        \node[zonelbl, anchor=south west] at (zC.north west) {Camada de dados e ingestão};
    \end{tikzpicture}%
    }
    \caption{Arquitetura e fluxo operacional do SIM na instância piloto de Caeté.}
    \label{fig:arquitetura}
\end{figure}
